\chapter{Hello, maya}
\label{ch:intro}

This book teaches you \maya{} --- a library for building programs that
draw their interface inside a terminal window --- assuming you are new
to almost everything involved. We will not skip steps. By the end of
this first chapter you will know what kind of program \maya{} is, why
anyone bothers building such a thing in 2025, what is happening on
your screen \emph{right now} as you read this in a terminal, and how
to compile and run your first \maya{} program. None of that requires
prior knowledge of \cpp{}; we will treat the example as a thing to
look at, and Chapter~\ref{ch:cpp} will teach you the language behind
it.

\section{What is a terminal? (the honest answer)}

You have almost certainly used a terminal: it is the black-or-dark
window that says things like \code{\$} or \code{>} and lets you type
commands. It might be called Terminal, iTerm, kitty, Alacritty, GNOME
Terminal, Konsole, Windows Terminal, or just \emph{the console}.

Here is what the terminal actually is, in plain words:

\begin{enumerate}[leftmargin=*]
  \item A \textbf{window on your screen} that shows a rectangular grid
    of text. Every position in the grid holds \emph{one character} ---
    a letter, a digit, a symbol --- possibly with a colour.
  \item A \textbf{program} (the \emph{terminal emulator}) that owns
    that window. The terminal emulator's job is twofold:
    \begin{itemize}[leftmargin=*]
      \item When the user types on the keyboard, send those keystrokes
        to whatever program is running inside.
      \item When that program prints text, draw the text into the
        window.
    \end{itemize}
  \item A \textbf{program running inside it}. Usually this is a
    \emph{shell} like \code{bash} or \code{zsh} --- the thing that
    shows the \code{\$} prompt --- but it can be \emph{any} program.
    When you run \code{ls}, the shell starts \code{ls}, which prints
    file names, then exits, and you are back at the shell.
\end{enumerate}

So the terminal emulator is a kind of two-way pipe with a screen
attached: text goes in (from the keyboard), text comes out (from the
program), and the rectangle of cells on screen is updated to reflect
what the program has printed so far.

\begin{idea}
A terminal is a window that displays a grid of characters, and a
program that pumps text between your keyboard and whatever else is
running.
\end{idea}

\subsection{Why does it look so old?}

Because it is. The protocol the terminal speaks --- the rules for
``move the cursor to row 5, column 12'' and ``print this in red'' ---
was finalised on a physical device called the DEC VT100 in 1978, and
every modern terminal still understands it. Programs written today
talk the same byte language as programs written for an actual VT100.
This is unusual in computing and turns out to be very valuable: any
program that follows the rules will work on every machine you can
\code{ssh} into, today, forever.

\section{What is a TUI?}

A \textbf{TUI} is a \emph{Text User Interface}: a program that does
not just dump output line by line, but \emph{takes over} the terminal
window and draws a proper interface inside it. Menus, panels,
scrollable lists, progress bars, even charts --- all rendered out of
characters.

You have used TUIs without thinking of them as such:

\begin{itemize}[leftmargin=*]
  \item \code{vim} or \code{nano} when you edit a file.
  \item \code{htop} when you watch processes.
  \item \code{git log} (the pager part), \code{less}, \code{man}.
  \item Installer screens on a Linux ISO.
\end{itemize}

A \textbf{GUI} (Graphical User Interface) program asks the operating
system to paint individual pixels: \emph{draw a blue rectangle from
$(100, 80)$ to $(300, 120)$}. A TUI program does something more
modest: it asks the terminal emulator to put specific characters into
specific grid cells, possibly with colour. The terminal emulator
itself decides how each character looks (the font, the cell size); the
TUI program only chooses what character goes where.

\begin{idea}
GUI = pixels you place anywhere. TUI = characters you place inside a
fixed grid of cells.
\end{idea}

That restriction --- one character per cell --- is the entire reason
TUIs are smaller, faster, and easier to write than GUIs. There is no
font engine, no anti-aliasing, no input method editor; just a grid and
some bytes.

\subsection{Why bother with TUIs in 2025?}

You might reasonably ask: why not just write a web page? Several
reasons:

\begin{itemize}[leftmargin=*]
  \item \textbf{They run anywhere there is a terminal.} A developer
    on a server they reached by \code{ssh} has no browser, no
    desktop, sometimes not even a graphical session. They do have a
    terminal.
  \item \textbf{They are very fast.} A typical TUI redraw is a few
    kilobytes of bytes. A web page is megabytes of HTML, CSS, JS,
    images, and a browser.
  \item \textbf{They are scriptable.} A TUI program's output is
    still text, so it composes with pipes, \code{grep}, \code{ssh},
    \code{tmux}, and the rest of the Unix toolbox.
  \item \textbf{They are pleasant to use for developers.} Keyboard
    first, no mouse round-trips, low latency, look good in any
    monospace font.
\end{itemize}

The trade-off is the visual restriction: no images (well, mostly), no
freeform layouts, no animations beyond what a grid of glyphs can do.
For a large class of programs, that is exactly the right trade.

\section{One frame of a TUI: what literally happens}
\label{sec:intro:one-frame}

To make the rest of this book concrete, let us walk what happens in
one ``frame'' of a typical TUI --- the unit of work between two
updates of the screen. We will use a counter app as the example: it
shows a number, and the keys \code{+} and \code{-} change it.

Suppose the screen currently shows the number $7$. The user presses
\code{+}.

\paragraph{1. The keyboard sends a byte to the kernel.}
When you press \code{+}, the OS kernel receives the byte $0x2B$
(ASCII for \code{+}) from your keyboard driver and places it into
the input buffer of whichever terminal program currently has focus
--- your terminal emulator.

\paragraph{2. The terminal emulator forwards the byte.}
Your terminal emulator (kitty, iTerm, GNOME Terminal, \dots) writes
the byte into one end of the \emph{pseudo-terminal} --- the kernel
byte-pipe between the emulator and the program running inside it.
Chapter~\ref{ch:terminal} covers the pseudo-terminal in detail.

\paragraph{3. The maya program reads the byte.}
The \maya{} program is sitting in a loop, asking the OS ``is there
any input ready?''. Now there is. It reads $0x2B$ from its standard
input.

\paragraph{4. Maya turns the byte into an Event.}
A single byte is not yet a useful concept. \maya{}'s input parser
decides ``this is a printable character, namely the character
\code{+}'' and constructs a structured value: \code{Key\{.kind =
KeyKind::Char, .character = '+'\}}. That value is then wrapped in
the broader \code{Event} type and handed to your application code.

\paragraph{5. Your code computes a new state.}
Your application receives the \code{Event}, decides it means
``increment'', and produces the new state: the counter is now $8$.
In an Elm-style \maya{} program this is one pure function:
\code{update(state=7, event=Inc) = 8}.

\paragraph{6. Maya turns the new state into an Element tree.}
Your \code{view} function is called with the new state ($8$) and
returns an \code{Element}: a description of the UI, ``a box with a
rounded border containing the text \code{Count: 8}''. This is a
value in memory, not yet bytes on the wire.

\paragraph{7. The layout engine assigns positions.}
Given the current terminal size --- say 80 columns by 24 rows ---
\maya{} walks the element tree and decides where every piece goes:
the outer box at $(0, 0)$ of size $14 \times 3$, the inner text at
$(1, 1)$, and so on. The output is a list of styled cells, one per
position.

\paragraph{8. The renderer builds a virtual screen.}
Those cells are written into a 2-D grid in memory --- the
\emph{canvas}. Every position on the screen now has a known
character, foreground colour, background colour, and attribute set.

\paragraph{9. The differ finds what changed.}
A copy of last frame's canvas is still around. \maya{} compares the
two canvases cell by cell. For our counter, almost everything is
identical; only one cell changed --- the \code{7} became an \code{8}.

\paragraph{10. Maya emits ANSI bytes for only that cell.}
A short string of bytes is produced: ``move cursor to (col, row)
of the changed cell; set the right colour; print \code{8}''. Even
though there might be hundreds of cells on screen, only a dozen
bytes are sent.

\paragraph{11. Those bytes flow back through the pty to the emulator.}
The emulator reads them, parses them, and updates its window. The
user sees the digit change. From keypress to pixel update is
typically a couple of milliseconds.

That is one frame. The rest of the book is, in a sense, just
detailed coverage of each of those eleven steps. The DSL
(Chapter~\ref{ch:dsl}) and Element tree
(Chapter~\ref{ch:elements}) cover steps 5--6; the layout engine
(Chapter~\ref{ch:layout}) is step 7; the canvas
(Chapter~\ref{ch:canvas}) is step 8; the differ
(Chapter~\ref{ch:simd}) is step 9; ANSI
(Chapter~\ref{ch:ansi}) is step 10; the runtime that orchestrates
the whole loop (Chapter~\ref{ch:elm}) is what sits around steps
3--6.

\begin{idea}
One frame: a byte arrives, it becomes an event, you compute new
state, maya builds the new screen, compares it to the old one, and
sends bytes for just the differences. That tight loop is what every
TUI library is, including \maya{}.
\end{idea}

\section{What is \maya{}?}

\maya{} is a \cpp{} \textbf{library} --- a bundle of code you include
in your own program --- whose job is to make writing a TUI pleasant.
Without a library, writing a TUI means manually printing escape
sequences (we will look at the bytes in Chapter~\ref{ch:terminal}),
manually tracking what is on screen, manually parsing keyboard input,
and manually re-laying-out everything when the window resizes.

With \maya{}, you describe what you \emph{want} on screen as a tree of
boxes and text, and \maya{} works out the bytes to send. Concretely,
the work the library does for you is:

\begin{enumerate}[leftmargin=*]
  \item Take your description (\emph{``a box containing the text
    `Hello' with a rounded border''}) and figure out the position
    and size of every piece given the current window size.
  \item Build a virtual grid of cells representing what the screen
    should look like.
  \item Compare that grid to what is \emph{already} on screen, and
    emit just enough escape sequences to update the cells that
    changed.
  \item Read keyboard, mouse, paste, and resize events and feed them
    back to your program.
  \item Restore the terminal to a clean state when your program
    exits, even on crashes.
\end{enumerate}

This book covers each of those jobs in detail. For now, treat
\maya{} as a black box that turns ``description of a UI'' into ``the
correct bytes on the wire''.

\section{What is a tree, and why a tree?}
\label{sec:intro:tree}

In computing, a \textbf{tree} is a data structure that looks like a
family tree turned upside down: one thing at the top, called the
\textbf{root}, which has zero or more \textbf{children}, each of
which may have its own children, and so on. The end of every branch
--- a node with no children --- is called a \textbf{leaf}.

A terminal UI is naturally tree-shaped, and once you see it you
cannot unsee it. Consider a chat application:

\begin{center}
\small
\begin{verbatim}
  +-------------------------------------------+
  | Title: my chat                           |
  +-------------------------------------------+
  | sidebar  |  message list                  |
  |  - alice |   alice: hi                    |
  |  - bob   |   me   : hello                 |
  |          |   bob  : evening               |
  +----------+--------------------------------+
  | type here...                             |
  +-------------------------------------------+
\end{verbatim}
\end{center}

Described as a tree, this UI is:

\begin{center}
\small
\begin{verbatim}
  Root: vertical stack
  |-- Title bar (text)
  |-- Body: horizontal split
  |   |-- Sidebar: vertical list
  |   |   |-- "alice"  (text)
  |   |   `-- "bob"    (text)
  |   `-- Messages: vertical list
  |       |-- "alice: hi"      (text)
  |       |-- "me   : hello"   (text)
  |       `-- "bob  : evening" (text)
  `-- Input: bordered text box
\end{verbatim}
\end{center}

Four kinds of node appear: ``vertical stack'', ``horizontal split'',
``list'', ``text''. The first three are \emph{container} nodes ---
they have children that they arrange. Text nodes are leaves: they
have no children, just a string.

A tree is the right structure because UIs are recursive: a sidebar
\emph{is} a UI; a message list \emph{is} a UI; both are pieces of a
bigger UI. Trees compose. You can reuse the message-list subtree in
a completely different program by grafting it onto a different root.

\maya{}'s \code{Element} type is exactly the tree representation of
a UI. Container nodes (\code{BoxElement}, \code{ElementList}) hold
\code{std::vector<Element>} of children; leaf nodes
(\code{TextElement}) hold the bytes of a string. The DSL is the
surface syntax that lets you spell that tree out as readable \cpp{}
code. Chapter~\ref{ch:elements} dissects the \code{Element} type in
full.

\begin{idea}
A UI is a tree of containers and text. Containers arrange their
children; text leaves carry the actual characters. \maya{}'s
\code{Element} type \emph{is} that tree.
\end{idea}

\section{Why does \maya{} exist when there are already others?}

There are many TUI libraries. \code{ncurses} (C, ancient and
respected), \code{ratatui} (Rust), \code{bubbletea} (Go),
\code{textual} (Python), \code{ink} (Node.js). Each is good in its
own way. \maya{} picks three positions that distinguish it.

\paragraph{1. Wrong UI code does not compile.}
Most libraries accept your code happily and let you discover the bug
when you run the program: a border colour set without a border becomes
a no-op, a negative padding silently becomes zero, a layout meant to
be a row accidentally becomes a column. \maya{} is built on \cpp{}'s
\emph{type system} so that a lot of these mistakes refuse to compile.
The compiler tells you, at build time, exactly which line is wrong.

As a tiny taste: in \maya{} you might write

\begin{lstlisting}
auto ui = t<"Hello"> | bcol<60, 65, 80>;   // border colour, no border yet
\end{lstlisting}

A library written in Python or Go would accept that and silently do
nothing with the border colour at run-time. \maya{} refuses to
compile it, because the \code{bcol} (``border colour'') modifier is
defined to only apply to nodes that have already been given a
border. The compiler error names the constraint that failed; you fix
the code by adding \code{| border\_<Round>} first.

You will not see why this is possible until Chapter~\ref{ch:dsl}; the
short version is that the description you write is stored in the
\emph{types} of the \cpp{} values, and \maya{}'s API states the rules
between those types.

\paragraph{2. Re-drawing is wired to be fast.}
Once \maya{} has decided what the next frame should look like, it
compares that virtual grid of cells with the grid that is currently on
screen. Only the cells that changed are sent to the terminal. That
comparison is done with the CPU's \textbf{SIMD} instructions ---
special instructions that operate on many bytes at once. The effect
is that even on a slow \code{ssh} link, a \maya{} program redraws
smoothly. We will dissect the SIMD code in Chapter~\ref{ch:simd}; you
do not need to know what SIMD is yet.

\paragraph{3. The runtime is borrowed from a language called Elm.}
\maya{} organises a whole program around one idea: your program's
state is a single value, your program's events are a single closed
list of cases, and the function that updates the state given an event
is \emph{pure} --- meaning it computes a new value and does nothing
else (no printing, no opening files, no network calls). Anything that
needs to ``happen in the world'' is described as data and handed to
the library to perform. This is unusual for \cpp{}, and the payoff is
programs that are easy to test, easy to replay, and easy to reason
about. Chapter~\ref{ch:elm} is dedicated to this idea.

If those three sentences mean nothing to you yet, that is fine. The
book is structured so each one is earned, slowly, before you have to
use it.

\section{The layers of a \maya{} program}

It will help to have a picture of how the work flows, even if many of
the words are unfamiliar now. Read it once; you do not need to
memorise it.

\begin{center}
\small
\begin{tabular}{r|l}
\textbf{Your code} & describe a UI; say what should happen on each event \\
\hline
\textbf{DSL} & \cpp{} syntax (\code{|}, \code{t<"...">}) that builds your description \\
\textbf{Element} & in-memory tree of boxes and text \\
\textbf{Layout} & gives every box a position and size \\
\textbf{Canvas} & a grid of cells with characters, colours, attributes \\
\textbf{Diff} & finds which cells changed since last frame \\
\textbf{ANSI} & escape-sequence bytes for the changed cells \\
\textbf{Writer} & sends those bytes to the terminal \\
\hline
\textbf{Terminal} & draws characters into the window \\
\end{tabular}
\end{center}

A few words on the unfamiliar ones:

\begin{itemize}[leftmargin=*]
  \item \textbf{DSL} stands for \emph{Domain-Specific Language}. It is
    not a separate language --- it is just normal \cpp{} ---  but the
    surface looks specialised. You write
    \code{t<"Hi"> | Bold | border\_<Round>} and that is valid \cpp{}
    that produces a UI description. The DSL is what makes \maya{}
    code short.
  \item \textbf{Element} is the name we give to the in-memory shape
    of your UI. It is a tree because boxes contain other boxes (and
    text). The element tree is the central data structure of
    \maya{}.
  \item \textbf{ANSI} is the byte language the terminal understands.
    \code{ANSI} is short for \emph{American National Standards
    Institute}, who blessed the standard in 1976. We will look at
    individual ANSI bytes in Chapter~\ref{ch:terminal}.
\end{itemize}

Each layer knows only about the one immediately below it. Your code
talks to the DSL, the DSL produces an Element tree, the layout engine
walks the tree, and so on. This separation is what makes the library
both testable and replaceable: someone could swap out the SIMD diff
for a slower scalar one without anything above noticing.

\section{Your first \maya{} program}

Here is a complete, real \maya{} program. It compiles, runs, prints a
small banner, and exits.

\begin{lstlisting}
#include <maya/maya.hpp>
using namespace maya;
using namespace maya::dsl;

int main() {
    constexpr auto ui = v(
        t<"Hello, maya"> | Bold | Fg<100, 180, 255>,
        t<"a tiny banner"> | Dim
    ) | border_<Round> | pad<1>;

    print(ui.build());
}
\end{lstlisting}

Before we decode it line by line, here is roughly what you should
see in your terminal after running this:

\begin{terminalbox}
\textcolor{mayablue}{\textbf{.----------------.}}\\
\textcolor{mayablue}{\textbf{|}}\, \textbf{Hello, maya}~~\, \textcolor{mayablue}{\textbf{|}}\\
\textcolor{mayablue}{\textbf{|}}\, a tiny banner \textcolor{mayablue}{\textbf{|}}\\
\textcolor{mayablue}{\textbf{`----------------'}}
\end{terminalbox}

A two-line text stack, the first line bold and pale blue, the second
dimmed, surrounded by a rounded border with one cell of padding
inside it. The corners are drawn with the Unicode box-drawing
glyphs the terminal renders as smooth round corners.

You are not expected to understand every token of the program. You
\emph{are} expected to be able to point at each piece and at least
name what it is. Let us go line by line.

\subsection{Line by line}

\paragraph{\code{\#include <maya/maya.hpp>}.}
\code{\#include} is the \cpp{} way of saying ``copy and paste the
contents of this file here before compiling''. \file{maya/maya.hpp} is
a single header that pulls in the public parts of \maya{}. After this
line, the names \code{v}, \code{t}, \code{Bold}, \code{print}, and so
on are available.

\paragraph{\code{using namespace maya;} and \code{using namespace maya::dsl;}.}
A \textbf{namespace} is a labelled container of names. \maya{}'s
public API lives inside the namespace \code{maya}, and the DSL pieces
live in \code{maya::dsl}. The two \code{using namespace} lines say
``for the rest of this file, names from those namespaces are visible
without their prefix''. Without them you would have to write
\code{maya::print(maya::dsl::v(...))}. Verbose; we elide it for
examples.

\paragraph{\code{int main()}.}
Every \cpp{} program starts at a function called \code{main}. The
\code{int} says \code{main} returns an integer (the exit status); we
do not write a \code{return} statement, and \cpp{} treats that as
\code{return 0} which means ``success''.

\paragraph{\code{constexpr auto ui = ...}.}
Three things to unpack:
\begin{itemize}[leftmargin=*]
  \item \code{auto} means ``figure out the type of this for me, based
    on the right-hand side''. \cpp{} normally requires you to spell
    the type; \code{auto} lets the compiler infer it. The actual type
    of \code{ui} here is something elaborate that we do not care to
    write out.
  \item \code{constexpr} means ``this value is computed at
    \emph{compile time}, not at runtime''. The whole tree --- the
    border, the colour, the text --- is baked into the binary. When
    the program starts, the tree is already there.
  \item \code{ui} is just the name of the variable. We could have
    called it anything.
\end{itemize}

\paragraph{\code{v(...)}.}
\code{v} is \maya{}'s function for ``stack these things vertically''.
The arguments inside the parentheses are the things to stack ---
in this case, two pieces of text. There is also \code{h(...)} for
horizontal.

\paragraph{\code{t<"Hello, maya">}.}
\code{t} is a \textbf{template} parameterised by a compile-time
string. The angle brackets \code{<\ldots>} are the template-argument
syntax. \code{t<"Hello, maya">} is a value representing the literal
text ``Hello, maya''. The peculiar thing here is that the text itself
lives inside the type --- swapping \code{"Hello"} for \code{"World"}
gives a different type. This is one of the tricks that lets the
compiler check things. Chapter~\ref{ch:cpp} explains how it works in
detail.

\paragraph{\code{| Bold | Fg<100, 180, 255>}.}
The vertical bar \code{|} is the \textbf{pipe operator} in \cpp{}. By
default it means bitwise-or on integers, but you can teach it to mean
anything for your own types. \maya{} has taught it to mean
\emph{compose this style onto the thing on the left}. \code{Bold} is
a tag value meaning ``apply bold''; \code{Fg<R, G, B>} is a template
that takes three numbers in the range $0$--$255$ and represents a
foreground colour.

So \code{t<"Hello, maya"> | Bold | Fg<100, 180, 255>} reads
left-to-right: ``the text `Hello, maya', made bold, given a foreground
colour of $(100, 180, 255)$''. A pleasant pale blue.

\paragraph{\code{| Dim}.}
Same idea: applies the \emph{dim} attribute (terminals usually render
this as a faded foreground). One of the few attributes whose
appearance varies most across terminals.

\paragraph{\code{| border\_<Round> | pad<1>}.}
After the inner \code{v(...)}, two more pipes wrap the result. The
first puts a rounded border around the whole stack; the second adds
one cell of padding (space) inside the border. \code{Round} is a
choice from a small set (\code{Round}, \code{Square}, \code{Double},
\code{Thick}, etc.); \code{pad<1>} carries the number $1$ in its type.

\paragraph{\code{print(ui.build())}.}
Finally, the call that does something. \code{ui.build()} converts the
compile-time description \code{ui} into a runtime \code{Element} ---
the in-memory tree we mentioned earlier. \code{print(...)} is the
simplest of four output modes \maya{} provides; it paints the element
once, to standard output, and returns. No event loop, no alternate
screen, no clearing the terminal. For a static banner, that is what
we want. Bigger programs use \code{run<P>(...)}, which we will meet in
Chapter~\ref{ch:elm}.

\subsection{What you should take away}

Three observations are worth carrying forward:

\begin{enumerate}[leftmargin=*]
  \item The structure of the UI lives in the \cpp{} source as a
    \emph{value}: \code{ui} is a thing you can pass to a function,
    return from a function, store in an array. It is not strings of
    markup, not a separate file format, not XML.
  \item Because the value is \code{constexpr}, the compiler builds
    it. If you wrote \code{Fg<300, 0, 0>} (out of range), the
    compiler would refuse. If you wrote \code{border\_<Squiggle>}
    (not a real style), the compiler would refuse. These checks are
    free at runtime --- they have already happened.
  \item Reading code with \code{|} is like reading English left to
    right: subject first, modifiers afterwards. Once you get used
    to it the syntax is quite pleasant to scan.
\end{enumerate}

\section{Building and running}

The \maya{} repository builds with \textbf{CMake}, a tool that
generates build files for your platform. You need a recent \cpp{}
compiler --- GCC 15 or newer, or a recent Clang --- because \maya{}
uses \cpp{}26 features. On Linux:

\begin{lstlisting}[style=shell]
git clone --recursive https://github.com/1ay1/agentty
cd agentty/maya
cmake -B build
cmake --build build -j$(nproc)
\end{lstlisting}

What each command does, briefly:

\begin{itemize}[leftmargin=*]
  \item \code{git clone --recursive ...} downloads the project and
    its sub-projects (called \emph{submodules}).
  \item \code{cd agentty/maya} moves into the maya directory.
  \item \code{cmake -B build} reads \maya{}'s \file{CMakeLists.txt}
    file and produces build files in a new directory called
    \file{build}.
  \item \code{cmake --build build -j\$(nproc)} actually compiles
    things, using as many parallel jobs as you have CPU cores.
\end{itemize}

If it succeeds, \file{build/} contains around two dozen example
programs --- \file{build/counter}, \file{build/dashboard},
\file{build/doom\_fire} --- one for each file in \file{examples/}.
Run any of them by typing its path:

\begin{lstlisting}[style=shell]
./build/counter
\end{lstlisting}

A small UI appears. Press the keys it documents (often \code{q} to
quit). Resize the terminal window while it runs; watch the UI
re-flow.

\subsection{Adding your own example}

The easiest way to try the program above is to drop it into
\file{maya/examples/hello.cpp}, re-run CMake, and rebuild:

\begin{lstlisting}[style=shell]
# from inside maya/
$EDITOR examples/hello.cpp        # paste the program
cmake -B build                    # picks up the new file
cmake --build build -j$(nproc)
./build/hello
\end{lstlisting}

This loop --- edit, build, run --- is the loop you will spend the
rest of the book in.

\begin{warning}
If the build complains about \code{constexpr}, \code{concept}, or
unknown standard-library headers, your compiler is too old. \code{g++
--version} should be 15 or newer. On Ubuntu/Debian:
\code{sudo apt install g++-15}. On macOS, install a recent Homebrew
\code{llvm}. On Windows, use WSL with a current Ubuntu.
\end{warning}

\section{What the rest of the book looks like}

The book has four parts.

\textbf{Part I, Foundations}, is the part you are reading. After this
chapter, Chapter~\ref{ch:terminal} explains what a terminal really is
at the byte level --- you will be able to write ``Hello'' in red
without any library at all by the end of it. Chapter~\ref{ch:cpp} is
a careful, from-the-ground-up tour of the \cpp{} features \maya{}
uses; if you are new to \cpp{}, this is the chapter that pays for
everything that follows.

\textbf{Part II, Architecture}, is about the way a \maya{} program is
structured. The Elm runtime (\ref{ch:elm}), the DSL that produces
UI descriptions (\ref{ch:dsl}), the in-memory tree that is the result
(\ref{ch:elements}), and the layout engine that decides where things
go (\ref{ch:layout}).

\textbf{Part III, Rendering}, is the bytes-on-the-wire half. Cells and
canvases (\ref{ch:canvas}), SIMD comparison (\ref{ch:simd}), the ANSI
language (\ref{ch:ansi}), and the special mode where \maya{} renders
\emph{inline} without taking over the terminal (\ref{ch:inline}).

\textbf{Part IV, Runtime and widgets}, is everything that happens
while a program runs. Fine-grained reactivity with signals
(\ref{ch:signals}), input parsing (\ref{ch:events}), the full effect
algebra (\ref{ch:cmdsub}), the widget catalogue (\ref{ch:widgets}),
and how to ship a program built on \maya{} (\ref{ch:build}).

You can read straight through. You can also skip Chapter~\ref{ch:cpp}
on first pass if you already know modern \cpp{} well, and come back
to it as a reference.

\section{How \maya{} compares to other TUI libraries}
\label{sec:intro:compare}

It helps to know what other people are doing. Here is the same
static banner --- text, bold, colour, border --- in four other
popular TUI libraries, alongside the \maya{} version, so you have a
felt sense of where this book is taking you.

\paragraph{ncurses (C, 1990s).}
\begin{lstlisting}[language=C]
#include <ncurses.h>
int main(void) {
    initscr();
    start_color();
    init_pair(1, COLOR_CYAN, COLOR_BLACK);
    box(stdscr, 0, 0);
    attron(A_BOLD | COLOR_PAIR(1));
    mvprintw(1, 2, "Hello, maya");
    attroff(A_BOLD | COLOR_PAIR(1));
    mvprintw(2, 2, "a tiny banner");
    refresh();
    getch();
    endwin();
}
\end{lstlisting}
Imperative. You position every character by row and column. There
is no layout; if the terminal is smaller than your hard-coded
numbers, the UI breaks.

\paragraph{textual (Python).}
\begin{lstlisting}[language=Python]
from textual.app import App
from textual.widgets import Static

class Hello(App):
    def compose(self):
        yield Static("[bold cyan]Hello, maya[/]\n[dim]a tiny banner[/]",
                     classes="box")
    CSS = ".box { border: round cyan; padding: 1; }"

Hello().run()
\end{lstlisting}
Declarative, pleasant, but markup is strings (the bold/colour
spans), and the CSS is a separate sub-language the Python compiler
cannot check.

\paragraph{ratatui (Rust).}
\begin{lstlisting}[language=C]
let block = Block::default()
    .borders(Borders::ALL)
    .border_type(BorderType::Rounded)
    .border_style(Style::default().fg(Color::Cyan));
let text = Text::from(vec![
    Line::from(Span::styled("Hello, maya",
        Style::default().fg(Color::Cyan).add_modifier(Modifier::BOLD))),
    Line::from(Span::styled("a tiny banner",
        Style::default().add_modifier(Modifier::DIM))),
]);
frame.render_widget(Paragraph::new(text).block(block), frame.size());
\end{lstlisting}
Close to \maya{} in spirit --- a tree of widgets, styles as values
--- but more verbose, and styles are runtime values, so e.g. an
out-of-range colour is caught only at runtime (or in your tests).

\paragraph{\maya{} (\cpp{}26).}
\begin{lstlisting}
constexpr auto ui = v(
    t<"Hello, maya"> | Bold | Fg<100, 180, 255>,
    t<"a tiny banner"> | Dim
) | border_<Round> | pad<1>;
print(ui.build());
\end{lstlisting}
The whole description is a \code{constexpr} value. The compiler
builds the tree at translation time; bad styles, bad borders, and
out-of-range colours are caught before the program runs.

The trade is roughly: more verbose than \maya{} in dynamic
languages; more strictly checked in \maya{} than anywhere else;
broadly similar mental model across the declarative tier (textual,
ratatui, \maya{}). Once you can read one, you can read all of them.

\section{What is a ``binary''?}
\label{sec:intro:binary}

We keep mentioning ``the binary'' --- the thing your compiler
produces. One sentence on what it is: a \textbf{binary} (or
\textbf{executable}) is a single file on disk full of CPU
instructions plus the data your program needs. When you run
\code{./hello}, the operating system loads that file into memory
and begins executing the instructions at a fixed entry point
(typically a function the compiler links in, which then calls your
\code{main}).

The binary is built once and run many times. Anything that can be
computed at compile time can live inside the binary directly ---
text strings, constant tables, style values --- without any cost at
startup. That is what \maya{} exploits with \code{constexpr}: the
UI tree of a static banner is, literally, encoded into the bytes of
the executable file.

\section{How this book is laid out}

A few visual conventions you will see in the chapters:

\begin{idea}
Blue boxes are \textbf{ideas} --- a single sentence summarising the
key intuition. If you only remember one thing from a section, this is
what it should be.
\end{idea}

\begin{theory}
Purple boxes are \textbf{theory} --- the background concept from
computer science or programming-language theory behind what the
section just explained. Skippable on first read; useful on second.
\end{theory}

\begin{warning}
Amber boxes are \textbf{pitfalls} --- a mistake people new to the
material commonly make, and how to avoid it.
\end{warning}

\begin{recap}
Green boxes at the end of a section summarise it in two or three
sentences.
\end{recap}

\begin{terminalbox}
Black boxes contain literal terminal bytes or output, the way the
terminal sees or shows them.
\end{terminalbox}

Code in monospace with a coloured rule on the left is \cpp{} (or
occasionally shell). Lines are numbered. Long examples are broken into
pieces and explained.

\begin{recap}
A TUI is a program that draws into the grid of cells inside a terminal
window. \maya{} is a \cpp{}26 library that lets you describe a TUI as
a tree of styled boxes and text and renders it for you. Its three
distinctive ideas --- compile-time correctness, SIMD-fast diffing, and
a pure-functional runtime --- are each given a chapter of their own
later. The first \maya{} program is a five-line static banner; the
loop you spent the rest of the book in is edit, build, run.
\end{recap}

\section*{Workshop}

You should not yet expect to do these without looking things up. They
are calibrated so that doing them takes you exactly through the
machinery the chapter introduced.

\begin{enumerate}[leftmargin=*]
  \item \textbf{Build the repository.} Get a clean
    \code{cmake --build build -j\$(nproc)} succeeding on your
    machine. If the build fails, the error message is the assignment
    --- read it, look up the missing tool, install it.
  \item \textbf{Run three examples} from \file{build/}. Suggested:
    \file{counter} (smallest), \file{dashboard} (mid-sized),
    \file{doom\_fire} (visually striking). Resize the window while
    each is running.
  \item \textbf{Compile the Hello-banner.} Drop the program from
    earlier into \file{maya/examples/hello.cpp}, re-run CMake,
    rebuild, run.
  \item \textbf{Mutate the banner.} Change the foreground colour to
    something else. Swap \code{Round} for \code{Double} and then
    \code{Thick}. Change \code{pad<1>} to \code{pad<3>}. Each
    change is one number or one identifier; each requires a rebuild
    to take effect.
  \item \textbf{Provoke a compile error.} Change \code{Fg<100, 180,
    255>} to \code{Fg<100, 180, 999>}. Read the compiler's error
    message. You do not need to fix it elegantly --- just observe
    that it caught the mistake.
\end{enumerate}

\section*{Glossary of terms introduced in this chapter}

Keep this list handy; every term here is referenced freely in the
rest of the book.

\begin{description}[leftmargin=*, style=nextline]
  \item[\textbf{ANSI}] The byte language a program writes to a
    terminal to control colour, cursor position, and screen modes.
    Standardised in 1976. Chapter~\ref{ch:terminal} covers it.
  \item[\textbf{Binary, executable}] A file on disk holding compiled
    CPU instructions for your program. \code{./hello} runs the
    binary called \file{hello}.
  \item[\textbf{Canvas}] \maya{}'s in-memory 2-D grid of styled
    cells; the working representation of the next screen.
    Chapter~\ref{ch:canvas}.
  \item[\textbf{Cell}] One position in the terminal grid. Holds one
    character (sometimes two-cell-wide), foreground colour,
    background colour, and attribute bits.
  \item[\textbf{\code{constexpr}}] \cpp{} keyword meaning ``compute
    this at compile time''. \maya{}'s DSL leans on it heavily.
  \item[\textbf{Diff}] The pass that compares the new canvas to the
    previous one and decides which cells changed.
    Chapter~\ref{ch:simd}.
  \item[\textbf{DSL}] Domain-Specific Language. The pleasant
    \code{t<"..."> | Bold | ...} syntax in \maya{} is a DSL
    embedded inside \cpp{}. Chapter~\ref{ch:dsl}.
  \item[\textbf{Element}] \maya{}'s in-memory UI tree. The output of
    your \code{view} function; the input to the layout engine.
    Chapter~\ref{ch:elements}.
  \item[\textbf{Elm runtime}] The Elm-inspired loop
    (\code{init}--\code{update}--\code{view}--\code{subscribe}) that
    \maya{} programs are organised around. Chapter~\ref{ch:elm}.
  \item[\textbf{Escape sequence}] A short string of bytes beginning
    with \code{ESC} (\code{0x1B}) that tells the terminal to do
    something. The grammar of ANSI.
  \item[\textbf{Frame}] One complete update cycle: read events,
    compute new state, build new canvas, diff, emit bytes.
  \item[\textbf{GUI}] Graphical User Interface --- pixel-based.
    Contrast with TUI.
  \item[\textbf{Layout}] The step that assigns each Element a
    position and size given the terminal dimensions.
    Chapter~\ref{ch:layout}.
  \item[\textbf{Leaf, node, root}] Tree vocabulary. The root is the
    top of a tree; leaves are nodes with no children; everything
    else is an internal node.
  \item[\textbf{NTTP}] Non-Type Template Parameter --- a template
    argument that is a value, not a type. The trick that lets
    \code{t<"Hello">} encode a string into a type.
    Chapter~\ref{ch:cpp}.
  \item[\textbf{Pseudo-terminal, pty}] The kernel-managed byte pipe
    between a terminal emulator and the program running inside it.
    Chapter~\ref{ch:terminal}.
  \item[\textbf{Raw mode}] A configuration of the terminal where the
    kernel passes every keystroke through immediately, with no line
    editing or echo. The mode every TUI runs in.
  \item[\textbf{SIMD}] Single Instruction Multiple Data: CPU
    instructions that operate on many bytes at once. \maya{}'s diff
    pass uses SIMD. Chapter~\ref{ch:simd}.
  \item[\textbf{stdin / stdout / stderr}] The three default file
    descriptors --- 0, 1, 2 --- a process inherits. In a TUI all
    three are wired to the pty.
  \item[\textbf{Tree}] A data structure: a root with children, each
    of which may have children. A UI is naturally a tree.
  \item[\textbf{TUI}] Text User Interface. A program that draws
    into the grid of cells inside a terminal window.
  \item[\textbf{UTF-8}] The standard variable-length encoding of
    Unicode characters into bytes. Modern terminals expect UTF-8.
\end{description}
